\documentclass[a4paper]{article}

% Basic packages
% Español (México) con XeLaTeX: fontspec para fuentes Unicode, polyglossia
% para la división silábica y los nombres del documento en la variante mexicana.
\usepackage{fontspec}
\usepackage{polyglossia}
\setdefaultlanguage[variant=mexican]{spanish}
% Los nombres chinos van como «pinyin (original)»: xeCJK da a los caracteres
% Han su propia fuente; sin ella XeLaTeX los omite con un aviso.
% FandolSong no cubre algunos caracteres de nombres (U+73FA); WenQuanYi Zen Hei sí.
\usepackage[AutoFallBack=true]{xeCJK}
\setCJKmainfont{FandolSong-Regular.otf}
\setCJKfallbackfamilyfont{\CJKrmdefault}{WenQuanYi Zen Hei}
% polyglossia no traduce los nombres de listings.
\AtBeginDocument{\providecommand{\lstlistingname}{}\renewcommand{\lstlistingname}{Listado}%
  \providecommand{\lstlistlistingname}{}\renewcommand{\lstlistlistingname}{Índice de listados}}
\usepackage{amsmath, amssymb}
\usepackage{graphicx}
\usepackage[margin=2.5cm]{geometry}
\usepackage[most]{tcolorbox}
\usepackage{etoolbox}
\usepackage{listings}
\usepackage{hyperref}
\usepackage{booktabs}
\usepackage{subcaption}
\usepackage{float}

% Highlight boxes
\newtcolorbox{knowledgebox}[1]{
    enhanced, colback=blue!5!white, colframe=blue!75!black, colbacktitle=blue!75!black,
    coltitle=white, fonttitle=\bfseries, title={#1}, attach boxed title to top left={yshift=-2mm, xshift=2mm},
    boxrule=1pt, sharp corners
}

\newtcolorbox{importantbox}[1]{
    enhanced, colback=yellow!10!white, colframe=yellow!80!black, colbacktitle=yellow!80!black,
    coltitle=black, fonttitle=\bfseries, title={#1}, sharp corners
}

\newtcolorbox{warningbox}[1]{
    enhanced, colback=red!5!white, colframe=red!75!black, colbacktitle=red!75!black,
    coltitle=white, fonttitle=\bfseries, title={#1}, sharp corners
}

% Listings
\lstset{
    language=Python,
    basicstyle=\ttfamily\small,
    keywordstyle=\color{blue},
    stringstyle=\color{red!60!black},
    commentstyle=\color{green!60!black},
    breaklines=true,
    frame=single,
    numbers=left,
    numberstyle=\tiny\color{gray},
    captionpos=b,
    extendedchars=true
}

% Front-page metadata
\newcommand{\notetitle}{Visual Prompting multimodal: SAM 1/2/3 y Grounding DINO}
\newcommand{\noteauthors}{Elaboradas a partir de materiales públicos del curso}
\newcommand{\notedate}{2025}
\newcommand{\videochannel}{Wudaokou Nash (五道口纳什)}
\newcommand{\videopublishdate}{2025}
\newcommand{\videoduration}{aproximadamente 19 minutos}
\newcommand{\videourl}{}
\newcommand{\videocoverpath}{cover.jpg}
\newcommand{\repourl}{https://github.com/hqhq1025/ai-course-notes}


% Footer with repo info
\usepackage{fancyhdr}
\pagestyle{fancy}
\fancyhf{}
\fancyfoot[C]{\thepage}
\fancyfoot[R]{\footnotesize\color{gray}\href{\repourl}{AI Course Notes}}
\renewcommand{\headrulewidth}{0pt}
\renewcommand{\footrulewidth}{0pt}

\begin{document}

\begin{titlepage}
\centering
{\Large Notas del curso\par}
\vspace{1.2cm}
{\huge\bfseries \notetitle\par}
\vspace{0.8cm}
{\large \noteauthors\par}
\vspace{0.3cm}
{\large \notedate\par}
\vspace{1.2cm}

\ifdefempty{\videocoverpath}{
{\small Al generar las notas, indique la ruta local de la portada del video.\par}
}{
\includegraphics[width=0.82\textwidth,height=0.45\textheight,keepaspectratio]{\videocoverpath}\par
}

\vfill
\begin{tcolorbox}[width=0.9\textwidth, colback=black!2!white, colframe=black!60, sharp corners]
\textbf{Autor o canal del video}: \videochannel\par
\textbf{Fecha de publicación}: \videopublishdate\par
\textbf{Duración del video}: \videoduration\par
\ifdefempty{\videourl}{}{
\textbf{Enlace del video}: \href{\videourl}{\nolinkurl{\videourl}}\par
}
\par
\textbf{Página del proyecto}: \href{\repourl}{\nolinkurl{\repourl}}\par
\end{tcolorbox}
\end{titlepage}

\tableofcontents
\newpage

%% --- Inicio del contenido --- %%

\section{Visual Prompting: concepto y motivación}

\subsection{Qué es Visual Prompting}

Visual Prompting (indicación visual) es una técnica que, en la interacción con modelos de visión y lenguaje (VLM) multimodales, superpone marcas visuales sobre la imagen original (como números, flechas, regiones enmascaradas, cajas delimitadoras, etc.) para ayudar al modelo a razonar y a ubicar objetos con mayor precisión.

\begin{importantbox}{La idea central de Visual Prompting}
La esencia de Visual Prompting es: antes de entregar la imagen al VLM, se usa un módulo de visión por computadora para preprocesarla (segmentación, detección, etiquetado con números), y luego se envía la imagen ya enriquecida al VLM para el razonamiento. De este modo, el problema de regresión de «generar coordenadas directamente» se convierte en el problema de clasificación de «elegir un número», lo que reduce en gran medida la dificultad del razonamiento del modelo.
\end{importantbox}

Tomando Set-of-Mark (SoM) como ejemplo: dada una imagen original y una pregunta en texto, SoM primero aplica un modelo de segmentación para hacer segmentación de instancias sobre toda la imagen, luego coloca una etiqueta numérica en el centro de cada región segmentada, y finalmente entrega esta nueva imagen con números, junto con la pregunta en texto, al VLM para que responda.

\begin{knowledgebox}{Por qué funciona Visual Prompting}
Modelos VLM como GPT-4V incluyeron, en la etapa de preentrenamiento, una gran cantidad de pares de imagen y texto provenientes de figuras etiquetadas (labelled figures) e ilustraciones de libros de texto (charts, textbooks). Las imágenes de estos datos de entrenamiento ya contienen de por sí flechas, números y regiones segmentadas como marcas visuales. Por eso, cuando agregamos artificialmente marcas similares sobre una imagen, esto activa de manera natural el mecanismo de atención interno del modelo hacia ese tipo de anotaciones, lo que mejora la precisión del razonamiento. Es similar a las anotaciones en los problemas de geometría matemática: números, flechas y marcas de dirección ayudan al modelo a concentrarse en la información clave.
\end{knowledgebox}

\subsection{De problema de regresión a problema de clasificación}

La ventaja más importante de Visual Prompting está en reducir la dificultad del razonamiento:

\begin{itemize}
    \item \textbf{Forma convencional}: el VLM necesita generar directamente las coordenadas exactas del objeto objetivo $(x, y)$ o la caja delimitadora $(x_1, y_1, x_2, y_2)$, lo cual es un problema de regresión
    \item \textbf{Forma con Visual Prompting}: la imagen ya está segmentada y numerada, así que el VLM solo necesita elegir el número correcto, lo cual se convierte en un problema de clasificación o selección
\end{itemize}

\begin{warningbox}{Visual Prompting no es una solución universal}
A medida que la capacidad visual de los VLM modernos (como GPT-4o, Gemini 2.5, etc.) se vuelve cada vez más fuerte, en algunos escenarios el modelo ya puede generar coordenadas exactas de forma directa. Pero en escenarios complejos (múltiples objetos, oclusiones, descripciones de posición, etc.), Visual Prompting sigue siendo una técnica auxiliar muy práctica. Al ayudar al modelo a concentrar su atención y a estabilizar el proceso de razonamiento, mejora de forma notable la precisión de la ubicación.
\end{warningbox}

\subsection{Resumen de la sección}

Visual Prompting es una técnica que superpone sobre la imagen original los resultados de un preprocesamiento de visión por computadora (máscaras de segmentación, números, cajas delimitadoras, etc.), lo que convierte el problema de regresión de la ubicación de objetos en un problema de selección, reduce la dificultad de razonamiento del VLM y mejora la precisión. Su eficacia proviene de que los datos de entrenamiento del VLM ya contienen de forma natural una gran cantidad de imágenes con anotaciones.

\section{La familia SAM: de la segmentación de instancias a la segmentación semántica}

\subsection{SAM 1 y SAM 2: segmentación a nivel de instancia}

SAM (Segment Anything Model) es la familia de modelos de segmentación de propósito general presentada por Meta. Tanto SAM 1 como SAM 2 admiten la segmentación a nivel de instancia, con las siguientes formas principales de entrada:

\begin{itemize}
    \item \textbf{Point Prompt}: se indica un punto en la imagen y se segmenta el objeto en el que se ubica dicho punto (por ejemplo, un punto en el centro del cuerpo de una ballena segmenta la ballena completa)
    \item \textbf{Bounding Box Prompt}: se da un cuadro delimitador y se segmenta el objeto dentro del cuadro
    \item \textbf{Automatic Mask Generation}: sin necesidad de ningún prompt de entrada, segmenta automáticamente todas las instancias de la imagen completa
\end{itemize}

\begin{knowledgebox}{Los prompts de Point/Box de SAM ya son, en sí mismos, Visual Prompting}
En sentido amplio, los prompts de Point y Bounding Box que acepta SAM ya son, en sí mismos, una forma de Visual Prompting: se indica una posición visual concreta en la imagen original para guiar la segmentación del modelo. Este diseño permite combinar SAM con otros modelos de manera flexible.
\end{knowledgebox}

La función de segmentación automática (Automatic Segmentation) de SAM 1 y SAM 2 no requiere ninguna entrada manual: basta con recibir la imagen original para segmentar todas las instancias de la imagen completa. Esta función es un componente fundamental de métodos de Visual Prompting como Set-of-Mark.

\subsection{SAM 3: segmentación semántica susceptible de prompt}

SAM 3 representa un salto cualitativo respecto a sus predecesores: admite la segmentación a nivel semántico y es promptable (susceptible de prompt):

\begin{importantbox}{Entrada multimodal de SAM 3}
SAM 3 admite las siguientes formas de entrada:
\begin{itemize}
    \item \textbf{Entrada de texto}: se ingresa una frase nominal simple (como ``cat''), y se segmentan todos los gatos de la imagen
    \item \textbf{Entrada Point/Box}: la forma de interacción clásica heredada de SAM 1/2
    \item \textbf{Entrada de imágenes de ejemplo}: se proporcionan imágenes de ejemplos positivos y negativos, y el modelo se acerca a los positivos y se aleja de los negativos
\end{itemize}
\end{importantbox}

\begin{warningbox}{Limitación de la entrada de texto de SAM 3}
La entrada de texto de SAM 3 solo admite frases nominales (noun phrase) simples; \textbf{no admite palabras de posición} (como «el gato de la izquierda» o «la taza de hasta arriba», descripciones que incluyen relaciones espaciales). Si se necesita procesar consultas complejas que incluyan descripciones de posición, hay que recurrir a un VLM para simplificar primero la descripción compleja a una frase nominal, y después entregarla a SAM 3.
\end{warningbox}

\subsection{Combinación de agente VLM + SAM 3}

Dado que SAM 3 no admite palabras de posición, se puede construir un Agent Pipeline de VLM + SAM 3 para procesar tareas complejas de localización de objetos:

\begin{enumerate}
    \item \textbf{El usuario ingresa una consulta compleja}: por ejemplo, «el niño más a la izquierda que trae puesto un chaleco azul»
    \item \textbf{El VLM simplifica la consulta}: convierte la descripción compleja en una frase nominal simple: «child in blue vest»
    \item \textbf{Segmentación con SAM 3}: a partir del texto simplificado, se segmentan todos los niños que traen puesto un chaleco azul
    \item \textbf{Asignación de números}: se asigna un número a cada resultado de segmentación
    \item \textbf{Segunda inferencia del VLM}: a partir de la imagen segmentada y numerada, el VLM puede ubicar con precisión al niño «más a la izquierda»
\end{enumerate}

La idea central de este Pipeline es: el VLM se encarga de la comprensión y el razonamiento del lenguaje, y SAM 3 se encarga de la segmentación visual; cada uno cumple su función y, en conjunto, completan tareas complejas de localización que un solo modelo difícilmente lograría.

\subsection{Resumen de la sección}

La familia SAM evolucionó de la segmentación de instancias de SAM 1/2 a la segmentación semántica susceptible de prompt de SAM 3. SAM 3 admite entradas multimodales de texto, punto, cuadro e imágenes de ejemplo, pero el texto se limita a frases nominales simples. Mediante la combinación de agente de VLM + SAM 3, se pueden procesar tareas complejas de localización que incluyen palabras de posición.
\section{Grounding DINO y el Pipeline de detección y segmentación}

\subsection{Grounding DINO: detección de objetos de vocabulario abierto}

Grounding DINO es un modelo de detección de objetos de vocabulario abierto (open-vocabulary) que localiza objetos en una imagen a partir de una descripción en lenguaje natural, y produce un cuadro delimitador (bounding box).

\begin{importantbox}{El Pipeline estándar de detección a segmentación}
En aplicaciones prácticas, la detección y la segmentación suelen usarse combinadas:
\begin{enumerate}
    \item Se usa Grounding DINO para detectar el objetivo a partir de una descripción textual y obtener el cuadro delimitador
    \item Se calcula el centro del cuadro delimitador
    \item Se ingresa el cuadro delimitador o el punto central como Prompt a SAM, para realizar una segmentación a nivel de píxel del objeto
\end{enumerate}
Así se refina la localización de grano grueso mediante un cuadro hacia una región de segmentación poligonal irregular a nivel de píxel.
\end{importantbox}

Como ejemplo, en un escenario de conducción autónoma: el usuario describe "un auto deportivo rojo", Grounding DINO primero enmarca con un cuadro delimitador el auto deportivo rojo en la imagen, y después SAM realiza sobre ese cuadro una segmentación precisa a nivel de píxel, obteniendo una región poligonal irregular.

\subsection{Grounding SAM 2: seguimiento de objetos a nivel de video}

Grounding SAM 2 combina Grounding DINO con la capacidad de seguimiento de video (tracking) de SAM 2, para el seguimiento de objetos a nivel de video:

\begin{enumerate}
    \item \textbf{Selección de un fotograma clave}: se toma un fotograma representativo del video
    \item \textbf{Detección del objetivo}: se usa Grounding DINO para detectar el objeto (por ejemplo, "little girl"), obteniendo el cuadro delimitador
    \item \textbf{Segmentación inicial}: se ingresa el cuadro delimitador como Box Prompt a SAM 2, segmentando el objetivo
    \item \textbf{Seguimiento en video}: se usa el módulo de Tracking de SAM 2 para seguir y segmentar continuamente ese objetivo en todos los fotogramas del video
\end{enumerate}

En la etapa de seguimiento se pueden usar varios tipos de información como prompt:
\begin{itemize}
    \item El cuadro delimitador original
    \item Puntos muestreados en la región de segmentación (Point Prompt)
    \item La región de Mask obtenida de la segmentación
\end{itemize}

\subsection{Resumen de la sección}

Grounding DINO ofrece capacidad de detección de vocabulario abierto, y combinado con la familia SAM forma el Pipeline estándar de detección a segmentación. Grounding SAM 2 extiende esto, además, al escenario de seguimiento en video, logrando un flujo de extremo a extremo de detección en un solo fotograma más seguimiento en todo el video.

\section{Set-of-Mark: numeración de toda la imagen para asistir la localización}

\subsection{Principio del método}

Set-of-Mark (SoM) es el trabajo representativo del Visual Prompting; el método es extremadamente simple pero muy efectivo:

\begin{enumerate}
    \item Se realiza segmentación de instancias sobre toda la imagen original (usando la función de segmentación automática de SAM)
    \item Se coloca un número en el punto central de cada región segmentada
    \item La nueva imagen con numeración, junto con la pregunta en texto, se introduce al VLM
    \item El VLM responde la pregunta de localización eligiendo un número
\end{enumerate}

\begin{importantbox}{Supuesto de diseño de SoM}
La efectividad de SoM se basa en un supuesto clave: durante el preentrenamiento, el VLM ya estuvo expuesto a una gran cantidad de labelled figures, charts y textbooks anotados. Estos datos contienen de manera natural marcas como números y flechas, por lo que el VLM tiene una fuerte capacidad de comprensión de este tipo de marcas visuales. SoM aprovecha esta capacidad: segmenta la imagen en regiones semánticas y asigna a cada región un ID legible, de modo que el VLM pueda «describir la imagen».
\end{importantbox}

\subsection{Análisis de ventajas}

Las dos ventajas centrales de SoM:

\begin{enumerate}
    \item \textbf{Reduce la dificultad de la inferencia}: convierte el problema de regresión de generar coordenadas directamente en un problema de clasificación de elegir un número
    \item \textbf{Ayuda a enfocar la atención}: las marcas numéricas tras la segmentación ayudan al modelo a enfocarse en regiones específicas, estabilizando la atención durante el proceso de inferencia
\end{enumerate}

\subsection{Resumen de la sección}

SoM, mediante la estrategia simple de segmentar toda la imagen y anotarla con números, convierte el problema de localización de objetos en una pregunta de opción múltiple, lo que reduce de manera notable la dificultad de inferencia del VLM, y constituye un ejemplo clásico de Visual Prompting.

\section{FreeGrasp: Visual Prompting en la sujeción robótica}

\subsection{Planteamiento del problema}

FreeGrasp resuelve el problema de la sujeción robótica de forma libre (free-form grasping): el usuario da una instrucción de sujeción ambigua en lenguaje natural (por ejemplo, «agarra el coche de juguete rojo que está abajo»), y el robot debe entender la instrucción, localizar el objeto objetivo, manejar las oclusiones y, por último, completar la sujeción.

\begin{knowledgebox}{El manejo de oclusiones es el reto central}
En escenarios reales de sujeción, el objeto objetivo suele quedar oculto por otros objetos. FreeGrasp diseñó un Pipeline completo para manejar las oclusiones: primero identifica el objetivo → determina si está oculto → retira el objeto que lo oculta → por último sujeta el objetivo. Este proceso requiere que el VLM tenga capacidad de razonamiento en varios pasos.
\end{knowledgebox}

\subsection{Diseño del Pipeline}

El flujo de procesamiento de FreeGrasp:

\begin{enumerate}
    \item \textbf{Entrada}: imagen RGBD (RGB para el razonamiento visual, Depth para la ejecución del brazo robótico) + instrucción en lenguaje natural
    \item \textbf{Detección de la imagen completa}: se usa un modelo de CV preentrenado para detectar todos los objetos de la imagen y obtener el punto central de cada uno
    \item \textbf{Visual Marking}: se marca un número en la posición de cada punto central
    \item \textbf{Razonamiento del VLM}: mediante un Prompt estructurado, se hace que el VLM razone en varios pasos a partir de la imagen con números:
    \begin{itemize}
        \item Identificar el objeto objetivo (¿qué número es el objetivo?)
        \item Determinar si está oculto (¿qué objetos ocultan al objetivo?)
        \item Planear el orden en que se retiran los objetos que lo ocultan
    \end{itemize}
    \item \textbf{Segmentación y sujeción}: se segmenta el objeto objetivo, un algoritmo tradicional estima la postura de sujeción y el brazo robótico la ejecuta
\end{enumerate}

\begin{importantbox}{El significado de Free-form}
En FreeGrasp, «Free» se refiere al grado de libertad de la entrada del usuario: se puede describir el objetivo de la sujeción con cualquier lenguaje natural, sin necesidad de un formato de instrucción estructurado. Esto se debe a la sólida capacidad de comprensión del lenguaje natural que tiene el VLM.
\end{importantbox}

\subsection{Resumen de la sección}

FreeGrasp aplica la técnica de Visual Prompting al ámbito de la sujeción robótica: mediante un Pipeline de detección → numeración → razonamiento del VLM → segmentación → ejecución, logra la sujeción de objetos a partir de instrucciones en lenguaje natural de forma libre, y la detección y el manejo de oclusiones son la innovación central.
\section{Pivot: Visual Prompting iterativo en la navegación robótica}

\subsection{Planteamiento del problema y método}

Pivot resuelve el problema de la navegación visual robótica: el robot necesita navegar y desplazarse en el entorno según una instrucción (por ejemplo, «encuentra la banca de madera»).

\begin{importantbox}{La optimización iterativa de Pivot}
Pivot utiliza una estrategia de Visual Prompting iterativa:
\begin{enumerate}
    \item Muestrea un conjunto de direcciones candidatas a partir del espacio de acciones
    \item Mapea las acciones candidatas al espacio de la imagen (marca flechas de dirección o números sobre la imagen)
    \item Hace que el VLM infiera, con base en la imagen marcada, las mejores direcciones candidatas
    \item Vuelve a muestrear dentro del subespacio seleccionado y repite el proceso anterior
    \item Tras varias iteraciones, produce la acción final
\end{enumerate}
\end{importantbox}

Las entradas centrales del algoritmo:
\begin{itemize}
    \item Imagen RGB original
    \item Instrucción de navegación
    \item Definición del espacio de acciones
    \item Número máximo de iteraciones $T$
    \item Cantidad de muestras por ronda $K$
\end{itemize}

En cada iteración, se muestrean $K$ acciones a partir de una política de acción probabilizada $\pi$; tras mapearlas al espacio de la imagen, se hace que el VLM elija los mejores candidatos, y luego se ajusta una nueva distribución de acciones; el ciclo continúa hasta alcanzar el número máximo de iteraciones. En esencia, es un proceso de optimización de muestreo, selección y refinamiento.

\subsection{El valor del Visual Prompting en la navegación}

En las tareas de navegación, el valor del Visual Prompting radica en convertir la selección abstracta de acciones (avanzar, girar 30 grados a la izquierda, girar 15 grados a la derecha, etc.) en una selección visual directa: se marcan sobre la imagen los números de las direcciones candidatas, y se hace que el VLM elija la dirección más razonable con base en la información visual.

\begin{warningbox}{La particularidad del escenario de navegación}
A diferencia de la localización de objetos, la navegación no requiere elegir «qué objeto de la imagen», sino decidir «hacia qué dirección moverse». Aquí, la función del Visual Prompting es discretizar el espacio continuo de direcciones en un número finito de opciones marcadas, lo que reduce la dificultad de decisión del VLM.
\end{warningbox}

\subsection{Resumen de la sección}

Pivot muestra la aplicación del Visual Prompting en la navegación robótica: mediante un flujo iterativo de muestreo, marcado, inferencia del VLM y refinamiento, logra decisiones de navegación basadas en información visual, y es un trabajo representativo de la extensión del Visual Prompting de la localización estática a la toma de decisiones dinámica.

\section{Síntesis y ampliación}

\subsection{El paradigma unificado del Visual Prompting}

Todos los trabajos presentados en esta entrega —SoM, FreeGrasp, Pivot, VLM+SAM Agent— siguen un paradigma unificado:

\begin{importantbox}{Paradigma unificado del Visual Prompting}
\begin{enumerate}
    \item Usar módulos clásicos de visión por computadora (segmentación, detección, muestreo, etc.) para preprocesar la imagen original
    \item Superponer marcadores visuales (numeración, flechas, máscaras, etc.) sobre el resultado del preprocesamiento
    \item Entregar la imagen enriquecida al VLM para el razonamiento y la toma de decisiones de alto nivel
\end{enumerate}
Esta combinación de «preprocesamiento con visión por computadora + razonamiento con VLM» aprovecha al máximo las ventajas complementarias de ambos.
\end{importantbox}

\subsection{Direcciones de evolución técnica}

\begin{itemize}
    \item \textbf{Mejora de la capacidad del modelo}: a medida que mejora la capacidad de comprensión visual de los VLM, algunos escenarios de Visual Prompting podrían dejar de ser necesarios, aunque conservan su valor en escenarios complejos
    \item \textbf{Integración de extremo a extremo}: en el futuro podrían surgir modelos unificados que integren de extremo a extremo el preprocesamiento con visión por computadora y el razonamiento con VLM
    \item \textbf{Más escenarios de aplicación}: desde la localización de objetos, la sujeción robótica y la navegación, hasta la conducción autónoma, las imágenes médicas y otros campos
    \item \textbf{Procesamiento dinámico de video}: combinado con la capacidad de seguimiento de video de SAM 2, el Visual Prompting puede extenderse a la comprensión y la interacción con video
\end{itemize}

\subsection{Lecturas adicionales}

\begin{itemize}
    \item Set-of-Mark Visual Prompting for GPT-4V (artículo original de SoM)
    \item Serie de artículos de Segment Anything Model (SAM 1/2/3)
    \item Grounding DINO: detección de objetos de vocabulario abierto
    \item FreeGrasp: sujeción robótica de forma libre
    \item Pivot: navegación con Visual Prompting iterativo
\end{itemize}

%% --- Fin del contenido --- %%

\end{document}
